\section{Per-character neuron clusters and the discrete-log lens}
\label{sec:character-clusters}

In \cref{sec:findings} we listed, as the first observation, that the
trained model converges on the multiplicative Fourier basis on
$\Zpstar{p}$ --- equivalent to additive Fourier on $\Zp{p-1}$ via the
discrete logarithm.  This note makes that statement visual.  For each
character $k \in \mathcal{K}$ we identify a cluster of MLP neurons
specialized to that character and reconstruct the cluster's contribution
to the residual stream as a function of the input pair $(a, b)$.  Two
companion images of that function appear together: one indexed by the
natural integer $a \in \{1, \ldots, p-1\}$, and one indexed by the
discrete logarithm $\dlog_g(a) \in \{0, 1, \ldots, p-2\}$.  In the
natural-integer indexing the cluster output looks like noise.  In the
discrete-log indexing it becomes clean diagonal stripes.  The
transformation between the two is what the embedding $\WE$ has learned to
implement.

We give the framework in \cref{subsec:cluster-framework},
present the canonical Legendre channel from \cref{sec:legendre-channel}
through the same lens in \cref{subsec:legendre-relooked}, then show three
character clusters from a new $d_{\text{mlp}} = 32$ model in
\cref{subsec:dmlp32-clusters}.  \Cref{subsec:dlog-explained} unpacks the
algebraic reason that dlog-sorting reveals periodicity, and what that
tells us about the network's representation.

\subsection{Framework}
\label{subsec:cluster-framework}

Let $\mathcal{K} \subset \widehat{(\Zpstar{p})}$ be the set of key
characters identified from the embedding spectrum (see
\cref{sec:findings}, item~3).  For each $k \in \mathcal{K}$ we identify
a cluster of MLP neurons as follows.  For each MLP neuron $j$ let
$v_j \in \R^p$ be its direct contribution to the c-axis of the logits,
\begin{equation}
  v_j \;=\; \WU^{\!\top}\,\Wout[\!:,\,j] \;\in\; \R^p.
  \label{eq:cluster-vj}
\end{equation}
Projecting $v_j$ onto the multiplicative-Fourier basis on $\Zpstar{p}$
gives the per-character output spectrum of neuron $j$.  Define
$\mathrm{dom}(j) \in \mathcal{K}$ as the character on which $|v_j|^2$ is
maximally concentrated, and the \emph{cluster for $k$} as
\begin{equation}
  \mathcal{C}_k \;=\; \{\,j \,:\, \mathrm{dom}(j) = k\,\}.
\end{equation}
We then compute the cluster's combined signal on the input grid by
running the model on all input pairs $(a, b) \in \Zpstar{p} \times
\Zpstar{p}$ and reading off, at the final position, the projection of
$\sum_{j \in \mathcal{C}_k} \Wout[\!:,\,j]\,h_j(a, b)$ onto the
unembed's $k$-character read direction.  Here $h_j(a, b)$ is the
post-ReLU activation of neuron $j$.  The reference signal we compare
against is the algebraic ideal
\begin{equation}
  R_k(a, b) \;=\; \cos\!\big[\theta_k(a) + \theta_k(b)\big],
  \qquad \theta_k(x) := \tfrac{2\pi k}{\,p-1\,}\,\dlog_g(x).
  \label{eq:reference-Rk}
\end{equation}
This is the real part of $\chi_k(a)\chi_k(b) = \chi_k(ab)$ and is what
character $k$'s MLP cluster ``should'' produce up to a global scaling.

\subsection{The Legendre channel, re-examined}
\label{subsec:legendre-relooked}

\Cref{sec:legendre-channel} introduced the order-$2$ Legendre channel
at $d_{\text{model}} = 24$: a $9$-neuron cluster that produces the
output $-\chi_{56}(ab) = -\chi_{56}(a)\chi_{56}(b)$ pointwise.
\Cref{fig:legendre-channel-relooked} reproduces the punchline figure
through the framework of \cref{subsec:cluster-framework} for
$k = 56$: in the natural-integer indexing (\emph{left}) the output is
a chaotic red/blue checkerboard, while sorted by discrete log
(\emph{right}) it cleanly resolves into the two QR/non-QR coset blocks
that exactly equal the algebraic reference $-\chi_{56}(ab)$.

\begin{figure}[ht]
  \centering
  \includegraphics[width=0.88\linewidth]{legendre_reconstructed_signal}
  \caption{The Legendre cluster at $d_{\text{model}} = 24$, viewed
  through the natural-vs-coset-sorted lens.  Top row is the model
  output; bottom row is the algebraic reference
  $-\chi_{56}(a)\chi_{56}(b)$.  The right column sorts rows and columns
  into QR/non-QR cosets, which for the order-$2$ character is the
  coarsest case of dlog sorting (just two equivalence classes).  The
  cluster output reproduces $-\chi_{56}(ab)$ pointwise.  Reproduced
  from \cref{fig:legendre-reconstructed} for direct comparison with the
  higher-order channels below.}
  \label{fig:legendre-channel-relooked}
\end{figure}

The Legendre case is the simplest case of the framework: an order-$2$
character partitions $\Zpstar{p}$ into just two cosets (QRs and
non-QRs), so the dlog-sorted view becomes a $2 \times 2$ block plot
rather than diagonal stripes.

\subsection{Three clusters from a $d_{\text{mlp}} = 32$ model}
\label{subsec:dmlp32-clusters}

To push the framework to higher-order characters we apply it to a
model trained with severe capacity restriction:
$d_{\text{model}} = 24$, $d_{\text{mlp}} = 32$.  This is the smallest
$d_{\text{mlp}}$ at which the model still groks (see
\cref{sec:findings}, item~10 and the d\_mlp sweep): test loss reaches
$\sim 10^{-4}$ at this width.  The trained model selects
$\mathcal{K} = \{3, 10, 20, 51\}$ with orders $\{112, 56, 28, 112\}$.
The neuron-output spectrum assigns all $32$ neurons to one of three
clusters $\mathcal{C}_3, \mathcal{C}_{10}, \mathcal{C}_{51}$
of sizes $11, 13, 8$ respectively.  No neuron's output is dominantly
on character $k = 20$.

\Cref{fig:cluster-k3,fig:cluster-k10,fig:cluster-k51} show the three
clusters' signals against their algebraic references.

\begin{figure}[ht]
  \centering
  \includegraphics[width=0.88\linewidth]{neuron_clusters_k3}
  \caption{The $11$-neuron cluster for $k = 3$ (order $112$) at
  $d_{\text{mlp}} = 32$.  Left column: natural-integer indexing of
  $(a, b)$.  Right column: rows and columns sorted by $\dlog_g$
  ($g = 3$ is a primitive root).  Top row: the cluster's combined
  signal projected onto the unembed's $k = 3$ read direction.  Bottom
  row: the algebraic reference $R_3(a, b)$ from
  \eqref{eq:reference-Rk}.  Correlation between model output and
  reference (over the whole $112 \times 112$ grid) is $+0.989$.  The
  dlog-sorted view is constant along anti-diagonals of constant
  $\dlog_g(a) + \dlog_g(b) \bmod (p-1)$, with cosine modulation at
  frequency $k = 3$ across diagonals.}
  \label{fig:cluster-k3}
\end{figure}

\begin{figure}[ht]
  \centering
  \includegraphics[width=0.88\linewidth]{neuron_clusters_k10}
  \caption{The $13$-neuron cluster for $k = 10$ (order $56$).
  Correlation with reference is $+0.965$.  The dlog-sorted view shows
  $112 / 10 \approx 11$ cosine oscillations along the anti-diagonal
  direction, twice as many as the $k = 3$ panel.}
  \label{fig:cluster-k10}
\end{figure}

\begin{figure}[ht]
  \centering
  \includegraphics[width=0.88\linewidth]{neuron_clusters_k51}
  \caption{The $8$-neuron cluster for $k = 51$ (order $112$).
  Correlation with reference drops to $+0.742$, reflecting the under-
  provisioning of this channel: $8$ neurons is at the edge of what
  suffices for an order-$112$ character.  The dlog-sorted view shows
  the correct stripe structure but with visible noise from incomplete
  implementation.  The model uses both $k = 3$ and $k = 51$ as
  primitive characters, with $k = 3$ comfortably implemented and
  $k = 51$ stretched thin.}
  \label{fig:cluster-k51}
\end{figure}

The three clusters together account for all $32$ MLP neurons, with
sizes $11 + 13 + 8 = 32$.  The model implements
\emph{three} character products at $d_{\text{mlp}} = 32$, despite
$\mathcal{K}$ having four elements: the fourth W$_E$-component
($k = 20$) does not have a dedicated MLP cluster, suggesting it is
either vestigial (kept alive by the embedding gradient but unused
downstream) or piggybacking weakly on the other clusters.

\subsection{Why dlog-sorting reveals periodic structure}
\label{subsec:dlog-explained}

The clean stripe pattern in the right column of each figure is not an
accident of plotting --- it is the visual signature of the
group-theoretic structure the network has learned.  Let $g$ be a
primitive root of $\Zpstar{p}$ and write $j_a := \dlog_g(a)$.  Every
element $a \in \Zpstar{p}$ has a unique exponent $j_a \in \Zp{p-1}$
such that $g^{j_a} \equiv a \pmod{p}$.

The character $\chi_k$ is defined by
\begin{equation}
  \chi_k(g^j) \;=\; e^{2\pi i k j / (p-1)},
\end{equation}
so $\chi_k$ is a group homomorphism from $\Zpstar{p}$ to the roots of
unity.  Crucially, $\chi_k$ is not a smooth function of $a$ --- but it
\emph{is} a smooth (indeed $C^\infty$) function of $j_a$, with period
$(p-1)/\gcd(k, p-1) = \ord(k)$.

The function plotted in each panel,
\begin{equation}
  R_k(a, b)
  \;=\; \cos\!\big[\theta_k(a) + \theta_k(b)\big]
  \;=\; \cos\!\left(\tfrac{2\pi k\,(j_a + j_b)}{p-1}\right),
\end{equation}
depends on $(a, b)$ \emph{only through the single number}
$s := j_a + j_b \bmod (p-1)$.  That is: in dlog coordinates the
$(p-1) \times (p-1)$ grid is reparameterized as a 1-D function pulled
back along the affine map $(j_a, j_b) \mapsto j_a + j_b$.  Two
consequences:

\begin{enumerate}
\item In dlog coordinates, $R_k$ is constant along every anti-diagonal
  $\{(j_a, j_b) : j_a + j_b = s\}$.  Anti-diagonals appear as straight
  lines of constant color, and the value cycles between them at
  frequency $k$ with wavelength $(p-1)/k$ (when $k$ is coprime to
  $p-1$) or $\ord(k)$ (in general).  For $k = 3$ the wavelength is
  $\sim 37$ steps along the diagonal direction, so $\sim 3$ full
  cosine periods are visible across the image; for $k = 10$ it is
  $11.2$, giving $\sim 10$ periods; for $k = 51$ it is $2.2$, giving
  fine high-frequency stripes.

\item In natural-integer coordinates, the rows and columns are
  permuted by the map $j \mapsto g^j \bmod p$, which is the discrete-
  log permutation in reverse.  This map has no efficient algebraic
  structure (its hardness is the security premise of discrete-log
  cryptography), so applying it to a smooth 1-D periodic function
  produces an apparently noisy 2-D pattern.  The structure is still
  there but encrypted by the choice of representation.
\end{enumerate}

The chaotic-looking left columns and the clean right columns are
therefore literally the \emph{same} function, displayed in two
different coordinate systems for $\Zpstar{p}$.

What this tells us about the network's representation is the punchline
of the section.  Modular multiplication has \emph{no} naturally smooth
structure in the integer representation; computing it via a lookup
table would require $\sim 12{,}544$ entries (for $p = 113$) and no
generalization across pairs $(a, b)$.  But $\Zpstar{p} \cong \Zp{p-1}$
via $\dlog_g$, and in the latter representation multiplication
\emph{becomes addition}.  Cosines of multiples of $j_a + j_b$ are then
smooth one-dimensional periodic functions, requiring only
$O(|\mathcal{K}|)$ characters to disambiguate the $p-1$ possible
outputs.  The network has discovered this isomorphism --- it has
effectively rediscovered the group structure --- and built its
algorithm in dlog coordinates.  The embedding $\WE$ implements the
isomorphism (taking $a \in \Zpstar{p}$ to a vector whose components are
$\cos\theta_k(a)$ and $\sin\theta_k(a)$, parameterized by $j_a$);
attention combines two such vectors additively in the $j$-coordinate;
the MLP applies the cosine product-to-sum identity to obtain
$\cos[\theta_k(j_a + j_b)]$; and the unembed reads off the result.

\subsection{Implications}
\label{subsec:cluster-implications}

\paragraph{Reconstruction quality scales with cluster size.}  Across
the three clusters at $d_{\text{mlp}} = 32$:
$11$ neurons $\to 0.989$ correlation;
$13 \to 0.965$;
$8 \to 0.742$.
The reconstruction degrades smoothly with neurons-per-character.  At
$d_{\text{model}} = 24$, $d_{\text{mlp}} = 512$ the same model
implements clusters of $\sim 50$--$100$ neurons per character with
near-perfect reconstruction; the cleanest minimal-circuit case is the
$9$-neuron Legendre cluster of \cref{sec:legendre-channel}, where the
$\chi_{56}$ output is exact within numerical noise.

\paragraph{The per-character $d_{\text{mlp}}$ floor.}  The empirical
floor for clean character reconstruction is around $\sim 8$--$13$
neurons per character.  Below that, channels degrade gracefully rather
than failing catastrophically: $k = 51$ at $8$ neurons is recognizable
($0.74$ correlation) but noisy.  This explains why the
$d_{\text{mlp}}$-cost lower bound of finding (5) in \cref{sec:findings}
($\sum_k \ord(k)$) is far too pessimistic --- the algorithm needs
$\sim 8$ neurons per character regardless of character order, not
$\ord(k)$ per character.

\paragraph{No new algorithm at tight $d_{\text{mlp}}$.}  Across
$d_{\text{mlp}} \in \{32, 64, 128, 256, 512\}$ the algorithm
implemented at the per-cluster level is the same Fourier-character
algorithm with no qualitative change.  Capacity pressure shrinks
clusters and reduces the number of characters in $\mathcal{K}$; it
does not produce a different computation strategy.  The basin geometry
(\cref{sec:basin-geometry}) commits to the same family of solutions
across all budgets that admit grokking.

\paragraph{Open question: what does the missing $k = 20$ do?}  At
$d_{\text{mlp}} = 32$ the embedding's top-energy character set is
$\{3, 10, 20, 51\}$, but only three clusters appear on the MLP side.
Two possibilities: ($i$) $k = 20$ is decoded by piggybacking on
neurons assigned to other clusters via small off-diagonal terms in
their output spectrum; ($ii$) $k = 20$ is a vestigial direction in
$\WE$ that does no algorithmic work and survives because the embedding
gradient (driven by the train cross-entropy) does not penalize it.
Distinguishing the two would require an ablation: remove the
$k = 20$ component from $\WE$ at test time and measure the change in
test loss.

\paragraph{Connection to discrete-log cryptography.}  The
``natural-integer chaos vs.\ dlog-sorted regularity'' duality in the
figures is exactly the same phenomenon that makes discrete-log
hard for adversaries.  Knowing the dlog map collapses modular
multiplication to addition; not knowing it forces a $p^2$-size lookup.
The trained network's embedding effectively contains the dlog
map (restricted to the characters in $\mathcal{K}$).  This raises a
question worth pursuing separately: under what circumstances does
training reveal information that is structurally hidden by the
cryptographic representation?  At $p = 113$ the answer is ``always
and trivially,'' because the dlog is computable by exhaustive search.
At cryptographic scales ($p \sim 2^{2048}$) it is presumed
intractable, but the same algorithmic content exists in the group ---
the network would just have to find it without enumerative help.
